\section{解析函数与 \texorpdfstring{$C^\infty$}{C-infinity} 函数}

\subsection*{1. 幂级数}

泰勒大胆地从插值公式得出了以他命名的级数展开式——泰勒级数展开式：
\begin{equation}
 f(x)=\sum_{k=0}^{\infty}\frac{f^{(k)}(a)}{k!}(x-a)^k.
 \tag{1}
\end{equation}
或者在多个变量即 $f:\Omega\subset\mathbb R^m\to\mathbb R^n$ 而 $m,n\ge1$ 的情况下有
\begin{equation}
 f(x)=\sum_{0\le|\alpha|<\infty}\frac{f^\alpha(0)}{\alpha!}(x-a)^\alpha.
 \tag{2}
\end{equation}
这里 $\alpha$ 是重指标。其实在那个时代什么是函数并没有划出清楚的界限。不但泰勒说不出他的结果是否适用于一切函数，甚至到了18世纪末，占主流的看法仍然认为所谓函数就是有限或无限的运算组合。而如(1)式这样的运算组合已经是够广泛的了。关于级数的收敛性，泰勒好像还没有这样一个概念。拉格朗日有一些初步的收敛性的概念，所以他是认为任意函数都可以写成(1)式。即令在今天，许多初学微积分的读者都会有一个错觉，以为只要 $f(x)$ 在 $a$ 附近为 $C^\infty$ 函数，则(1)式成立。这就大错特错了。一个古典的例子是考虑实变量 $x$ 在 $x=0$ 附近的函数
\[
 f(x)=
 \begin{cases}
 e^{-1/x},&x>0,\\
 0,&x\le0.
 \end{cases}
\]
用很简单的极限运算立即可知 $f(x)\in C^\infty$（包括 $x=0$），而且对一切非负整数 $k$，
\[
 f^{(k)}(0)=0.
\]
因此，对它应用(1)式，并取 $a=0$，则(1)式左方不恒为0而右方恒为0，可见(1)并非对一切 $C^\infty$ 函数都成立。因此，有必要将使得(1)式成立的函数专门划分为一类。这就是我们下面要讲的解析函数。

在给出解析函数定义以前，我们先要注意，(1)与(2)都是所谓幂级数——确切一些说是以 $x=a\in\mathbb R^m$ 为中心（center）的幂级数：
\[
 \sum_{n=0}^{\infty}a_n(x-a)^n
 \qquad\text{或}\qquad
 \sum_{0\le|\alpha|<\infty}a_\alpha(x-a)^\alpha.
\]
后式中 $\alpha$ 是重指标。这时有一个重要的说明：

1. 正如在第二章中讲到指数函数与三角函数时指出的，当一个函数可以用幂级数表示时，最好是立即进入复域。因此，在本节中讲到幂级数时，恒认为自变量 $x$，中心 $a$ 以及系数 $a_n$ 或 $a_\alpha$ 都是复数。

2. 在讨论自变量为复数的函数时，单复变量与多复变量的情况有重大的原则区别。因此，下面我们都只限于单自变量的函数 $f:\Omega\subset\mathbb C\to\mathbb C$ 的情况。$\mathbb C$ 指复平面，它具有复维数1，但是也可看作实维数2：$\mathbb C\simeq\mathbb R^2$。作为复自变量，我们总用 $z$ 表示，但 $z=x+iy$，$x,y$ 表示 $z$ 的实部与虚部，都是实数。同样复的函数值总用 $w=u+iv$ 表示。$u,v$ 是 $w$ 的实、虚部，都是实数。所以以下面讨论的复函数（通常此词是指复自变量函数，另一个词是复值函数，则是指函数值 $w$ 为复值，而自变量则不一定）有两种表示法：
\[
 w=f(z)
 \qquad\text{或}\qquad
 u=u(x,y),\ v=v(x,y).
\]
总之，本节中我们只讨论以下的幂级数：
\begin{equation}
 w=\sum_{n=0}^{\infty}a_nz^n
 \qquad\text{或}\qquad
 w=\sum_{n=0}^{\infty}a_n(z-a)^n.
 \tag{3}
\end{equation}
我们先从幂级数的一般性质的讨论开始，我们在第二章中已经说过，复变量函数的极限与连续性与复级数的收敛性、绝对收敛性及一致连续性等概念与基本的定理都与对应的实的情况是一样的。

幂级数的收敛性最基本的特点反映在下述阿贝尔定理中（我们只看(3)中 $a=0$ 的情况，一般情况相同）。

\noindent\textbf{定理1（阿贝尔（N. Abel）定理）}
对于级数(3)，必可找到一个非负实数 $R$（$R$ 允许取 $+\infty$ 值），使得

1）(3)在圆 $|z|<R$ 中绝对收敛，在 $|z|>R$ 处发散。

2）若 $K\subset\{z;|z|<R\}$ 是任一紧集，则(3)在 $K$ 上一致收敛。

圆 $|z|<R$ 称为(3)之收敛圆，$R$ 称为收敛半径。

\noindent\textbf{注}
定理中指出 $R\ge0$，若 $R=0$，则 $|z|<R$ 没有意义，而说(3)之收敛半径为0就是指对一切 $|z|>0$，级数(3)均为发散。

\noindent\textbf{证}
1）的证明。设(3)在一点 $z_0\ne0$ 处收敛，令 $\rho=|z_0|>0$。今证(3)在 $|z|<\rho$ 处必绝对收敛。事实上，因 $\sum_{n=0}^{\infty}a_nz_0^n$ 收敛，故其一般项必有界（实际上当 $n\to\infty$ 时趋于0）：$|a_nz_0^n|\le M$。考虑级数 $\sum|a_nz^n|$ 其中 $z$ 适合不等式 $|z|<|z_0|=\rho$，我们有
\[
 |a_nz^n|=|a_nz_0^n|\left|\frac z{z_0}\right|^n\le M\left|\frac z{z_0}\right|^n.
\]
但是 $|z/z_0|=|z|/\rho<1$，所以几何级数 $\sum_{n=0}^{\infty}M|z/z_0|^n<+\infty$，由比较判别法即知 $\sum_{n=0}^{\infty}|a_nz^n|<+\infty$，亦即(3)绝对收敛。

若(3)在一点 $z_0\ne0$ 处发散，则在 $|z|>|z_0|$ 处(3)亦必发散。事实上，若(3)在一点 $z_1$ 收敛，而 $|z_1|>|z_0|$，则把证明的前一部分用于 $\sum a_nz_1^n$，即知它是绝对收敛的，而与(3)在 $z_0$ 处发散相矛盾。

现在看一个集合 $J=\{\rho>0;(3)\text{在 }|z|<\rho\text{ 处绝对收敛}\}$。有两个可能，一是 $J=\varnothing$，即是说(3)当且仅当 $z=0$ 时收敛。这时我们取收敛半径 $R=0$。二是 $J$ 非空。先看 $J$ 非空的情况，令 $R=\sup J$，则 $R>0$。若 $z_0$ 点适合 $|z_0|<R$，现在要在 $J$ 中找一个数 $\rho$ 使 $|z_0|<\rho$（$R$ 还未证明是 $J$ 中之数）。记 $\varepsilon=R-|z_0|>0$，由于 $R$ 是 $J$ 之上确界，$J$ 中必有至少一个元 $\rho_0>R-\varepsilon/2$，而由定义知一定有：(3)在 $|z|<\rho_0$ 时绝对收敛。现在 $z_0$ 适合 $|z_0|=R-\varepsilon<R-\varepsilon/2<\rho_0$，所以(3)在 $z_0$ 处绝对收敛。即 $|z_0|<R$ 必导致(3)在 $z_0$ 处绝对收敛。所以实际上 $R$ 是 $J$ 中之数。又可以证明当 $|z_0|>R$ 时，(3)在 $z_0$ 处必发散。因为若(3)在 $z_0$ 处收敛，则令 $\varepsilon=|z_0|-R$，$\rho=|z_0|-\varepsilon/2=R+\varepsilon/2>R$，当 $|z|<\rho$ 时必有 $|z|<|z_0|-\varepsilon/2$，而(3)在 $z_0$ 处绝对收敛，这样 $\rho\in J$ 而与 $R=\sup J$ 矛盾。

总之，当 $J$ 非空时 $R$ 必是收敛半径而定理结论中的(1)成立。若 $J$ 为空，则取 $R=0$，结论(1)平凡地成立。

2）的证明如下。若 $R=0$，则 $K=\varnothing$ 而结论平凡地成立。若 $R>0$，则 $\{z;|z|<R\}$ 是一开集，其紧子集 $K$ 必有界的，而且 $K$ 必与 $\{z;|z|<R\}$ 之边界 $\{z;|z|=R\}$ 有正的距离，从而 $K$ 含于 $|z|\le R-\delta/2$ 中，且 $\sum_{n=0}^{\infty}|a_n|(R-\delta/2)^n$ 收敛。对于 $z\in K$，因为 $|a_nz^n|\le|a_n|(R-\delta/2)^n$，由一致收敛性的 $M$ 判别法知道 $\sum_{n=0}^{\infty}a_nz^n$ 在 $K$ 中一致收敛。

由阿贝尔定理得知幂级数一定有收敛半径 $R\ge0$。下面给出 $R$ 一个具体表达式。

\noindent\textbf{定理2（阿达玛）}
幂级数(3)的收敛半径是
\begin{equation}
 R=1\bigg/\overline{\lim}_{n\to\infty}\sqrt[n]{|a_n|}.
 \tag{4}
\end{equation}

\noindent\textbf{证}
证明分两部分，先证当 $|z|<R$ 时(3)必绝对收敛。首先，$\sqrt[n]{|a_n|}$ 的极限不一定存在，但上极限却是一定存在的，而且易见 $R\ge0$。如果 $R=0$，则 $|z|<R$ 是一个空集而谈不上(3)是否在 $z$ 点收敛或绝对收敛。如果 $R>0$，则由上极限的定义，对于任意小的正数 $\varepsilon>0$，最多除有限多个 $n$，一切系数均适合
\[
 \sqrt[n]{|a_n|}\le\frac1R+\varepsilon.
\]
现在 $z$ 是固定的，故当 $|z|<R$ 时，一定可以找到一个实数 $\lambda:0<\lambda<1$，使 $|z|\le\lambda R$。于是从某一个 $N_0$ 开始，当 $n>N_0$ 时
\[
 |a_nz^n|\le\left[\left(\frac1R+\varepsilon\right)\lambda R\right]^n.
\]
现在 $\lambda+1<2$，从而 $2/(\lambda+1)>1$，因此一定可以取 $\varepsilon>0$ 使得
\[
 \frac1R+\varepsilon<\frac2{(\lambda+1)R}.
\]
于是上式给出
\[
 |a_nz^n|\le\left(\frac{2\lambda}{\lambda+1}\right)^n.
\]
但是 $2\lambda/(\lambda+1)<1$，所以由级数绝对收敛性的比较法则（它对于复的无穷级数仍成立，这一点读者可以自己验证），即知(3)当 $|z|<R$ 时绝对收敛。

证明的另一部分是求证当 $|z|>R$ 时，级数(3)必发散。与上一部分不同，这一部分证明对 $R=0$ 也适用。由上极限的定义，对任一 $\varepsilon>0$，必有一个上升序列 $n_1<n_2<\cdots<n_k<\cdots\to+\infty$ 使
\[
 \sqrt[n_k]{|a_{n_k}|}\ge\frac1R-\varepsilon,
\]
所以
\[
 |a_{n_k}z^{n_k}|\ge\left[\left(\frac1R-\varepsilon\right)|z|\right]^{n_k}.
\]
但因 $z$ 是固定的且 $|z|>R$，所以必可写出 $|z|=R+\delta$ 而 $\delta$ 是一个固定的正数。适当取 $\varepsilon$ 使
\[
 \left(\frac1R-\varepsilon\right)(R+\delta)>1.
\]
代入上式即有 $|a_{n_k}z^{n_k}|\ge1$，即是说级数(3)的一般项不趋向0因而发散。注意若 $R=0$，上面用到 $1/R$ 处都没有意义了。但 $R=0$ 即 $\overline{\lim}\sqrt[n]{|a_n|}=+\infty$，因而对任意大的正数 $M$，必可找到子序列 $n_1<n_2<\cdots<n_k<\cdots$ 使 $\sqrt[n_k]{|a_{n_k}|}\ge M-\varepsilon$。上面凡用到 $1/R$ 处都换成 $M$，即知论证仍然有效，故定理成立。

凡讲到幂级数(3)之收敛性时，总有两个特例：一类幂级数收敛半径 $R=0$，即(3)只在中心 $z=0$ 收敛。例如 $\sum_{n=1}^{\infty}n!z^n$ 即是。但是我们不必用阿达玛公式(4)来计算 $R$，因为那会涉及求 $\overline{\lim}\sqrt[n]{n!}$，这个计算并不容易，而需要一个有名的、也非常有用的斯特林（Stirling）公式。第四章中我们还会遇到它。对于 $\sum_{n=0}^{\infty}n!z^n$ 反而用比值 $a_{n+1}/a_n$ 来求收敛半径要容易得多。另一类幂级数的收敛半径 $R=\infty$，即 $\overline{\lim}\sqrt[n]{|a_n|}=0$（也就是 $\lim\sqrt[n]{|a_n|}=0$，因为这时极限值一定存在）。于是级数(3)对一切 $z$ 均收敛，这时(3)之和称为整函数。数学中许多重要的函数如 $e^z,\cos z,\sin z,\ldots$ 都是整函数。不论如何，幂级数在收敛圆内的性状已完全清楚：在收敛圆外也很简单：发散。但是在收敛圆周上情况却十分复杂。(3)可能收敛，可能绝对收敛，可能一致收敛——一直到圆周……但也可能发散。总之，级数(3)在收敛圆周上的性状是一个专门的研究领域。

要注意，级数(3)在收敛圆内是绝对收敛的。无穷级数理论中有一个基本的事实：绝对收敛级数在代数运算性质上最接近于有限和：可以交换求和次序，适用分配律等等。因此，幂级数的运算性质最为简单。下面我们只举出几个最简单常用的性质而都不再证明。

首先是乘法。设有两个幂级数 $\sum_{n=0}^{\infty}a_nz^n$ 与 $\sum_{m=0}^{\infty}b_mz^m$，其收敛半径分别是 $R_1$ 与 $R_2$，则
\[
 \left(\sum_{n=0}^{\infty}a_nz^n\right)\left(\sum_{m=0}^{\infty}b_mz^m\right)
\]
本来是一个二重级数 $\sum_{m,n=0}^{\infty}a_mb_nz^{m+n}$。但是我们常用的是以下的写法，即按 $m+n=k$ 的大小来排列，即有
\begin{equation}
 \left(\sum_{n=0}^{\infty}a_nz^n\right)\left(\sum_{m=0}^{\infty}b_mz^m\right)
 =\sum_{k=0}^{\infty}\left(\sum_{m+n=k}a_mb_n\right)z^k.
 \tag{5}
\end{equation}
注意，内含的求和 $\sum_{m+n=k}a_mb_n$ 只是有限的，而不存在收敛性问题，故(5)式不是一个二重级数，(5)的收敛半径不小于 $\min(R_1,R_2)$。

其次是求函数的复合。设有两个函数 $F(w)=\sum_{n=0}^{\infty}a_nw^n$，其收敛半径为 $R_1$，以及 $\varphi(z)=\sum_{m=1}^{\infty}b_mz^m$，其收敛半径为 $R_2$，注意我们设 $b_0=0$，其理由从下面即可以看到。现在我们要求 $(F\circ\varphi)(z)$ 的幂级数展开式，一方面我们有
\begin{equation}
 F[\varphi(z)]=\sum_{n=0}^{\infty}a_n[\varphi(z)]^n.
 \tag{6}
\end{equation}
但是，为使它收敛，就需要 $|\varphi(z)|<R_1$。因为我们已设 $b_0=0$，所以当 $z=0$ 时 $\varphi(z)=0$，而当 $|z|$ 充分小时 $|\varphi(z)|$ 也充分小。我们假设 $b_0=0$ 的原因即在此。这里我们用到了 $\sum_{m=1}^{\infty}b_mz^m$ 在它的收敛圆 $|z|<R_2$ 的任一紧子集上一致收敛，因而其和在 $|z|<R_2$ 内连续。这里读者会问，为什么不说其和在 $|z|<R_2$ 之任一紧子集内连续？这是因为连续性是一点附近的性质，任取一点 $z_0\in\{z;|z|<R_2\}$，则 $z_0$ 必可放在一个闭圆 $K$ 中，而 $K\subset\{z;|z|<R_2\}$。$K$ 是一个有界闭集即紧集，而 $\sum_{m=1}^{\infty}b_mz^m$ 在 $K$ 上一致收敛，从而其和 $\varphi(z)$ 在 $K$ 上连续，当然也在 $z_0$ 连续。因为 $z_0$ 是收敛圆内任一点，所以 $\varphi(z)$ 在整个收敛圆内，而不只在其任一紧子集上连续。但是一致连续性并不是对于某一点来说的，而是要指定一个区域——现在是收敛圆，而我们恰好不能证明(3)在收敛圆内一致收敛。由 $\varphi(z)$ 在 $|z|<R_2$ 内连续以及 $\varphi(0)=0$ 即知一定有一个正数 $r>0$ 存在，使当 $|z|<r$ 时 $|\varphi(z)|<R_1$，从而级数(6)不但有有意义的各项，而且(6)收敛。于是
\begin{equation}
 F(\varphi(z))=\sum_{n=0}^{\infty}a_n\left(\sum_{m=1}^{\infty}b_mz^m\right)^n.
 \tag{7}
\end{equation}
但是我们可以利用上面讲的乘法来把 $\left(\sum_{m=1}^{\infty}b_mz^m\right)^n$ 写成幂级数
\begin{equation}
 \left(\sum_{m=1}^{\infty}b_mz^m\right)^n=\sum_{k=n}^{\infty}B_{nk}z^k.
 \tag{8}
\end{equation}
现在的问题是如何把它代入(7)式。为此要注意这个级数是由 $z^n$ 开始的，而且易见 $B_{nn}=b_1^n$。因此当我们把(8)代入(7)，并且要求例如 $z^l$ 之系数时，则(7)中 $n>l$ 的各项中一定不会有 $z^l$ 出现，$z^l$ 只出现在相应于 $n=0,1,\cdots,l$ 各项中，而成为
\[
 \sum_{n=0}^l a_nB_{nl}=A_l,\qquad\text{其中 }B_{00}=1.
\]
我们只知道它有一个正的收敛半径，但是没有办法把它用 $R_1,R_2$ 显式地表示出来。

以上我们只提供了一个计算复合函数的幂级数的程序，一般说来，它是相当繁冗的，但是它是相当重要的。许多很难的数学定理、数学公式都要靠这样的方法计算出来。下面我们只举一个例子，即求函数的倒数。

设级数(3)有收敛半径 $R$，其和为 $\varphi(z)$，我们设 $\varphi(0)=a_0\ne0$，现在来求 $1/\varphi(z)$，为此我们把 $\varphi(z)$ 写成
\[
 \varphi(z)=a_0+a_1z+\cdots=a_0\left[1-\sum_{m=1}^{\infty}b_mz^m\right],
\]
而 $b_m=-a_m/a_0$。于是
\[
 \frac1{\varphi(z)}=\frac1{a_0}\frac1{1-\sum_{m=1}^{\infty}b_mz^m}.
\]
由此可知，我们可以考虑
\begin{equation}
 F(w)=\frac1{a_0}\frac1{1-w}.
 \tag{9}
\end{equation}
再把它与 $w=\sum_{m=1}^{\infty}b_mz^m$ 复合起来即可。但是在整个“复分析”（即复变量函数的微积分学）中，“最重要”的幂级数展开式就是几何级数
\[
 \frac1{1-w}=1+w+w^2+\cdots=\sum_{n=0}^{\infty}w^n,
\]
它的收敛半径是1，这一点是所有读者在任何时刻都不可忘记的。应用它和上面所说的，即有
\[
 \frac1{\varphi(z)}=\frac1{a_0}\sum_{n=0}^{\infty}\left(\sum_{m=1}^{\infty}b_mz^m\right)^n,
\]
这里 $|z|$ 应充分小以保证 $\left|\sum_{m=1}^{\infty}b_mz^m\right|<1$。但是我们得不到收敛半径的具体表达式。

更具体一点，我们来看 $\sec z=1/\cos z$ 的幂级数展开式。现在因
\[
 \cos z=1-\frac{z^2}{2!}+\frac{z^4}{4!}\cdots,
\]
所以 $a_0=1$，而
\[
 \sum_{m=1}^{\infty}b_mz^m=\sum_{m=1}^{\infty}\frac{(-1)^{m-1}}{(2m)!}z^{2m},
\]
所以
\begin{align*}
 \sec z
 &=1+\left(\sum_{m=1}^{\infty}\frac{(-1)^{m-1}}{(2m)!}z^{2m}\right)
 +\left(\sum_{m=1}^{\infty}\frac{(-1)^{m-1}}{(2m)!}z^{2m}\right)^2+\cdots\\
 &=1+\frac12z^2+\frac5{24}z^4+\cdots.
\end{align*}
它的系数有些什么规律性，几乎无法看出。它的收敛半径是 $\pi/2$ 也是用其他方法得出的。总之，这是许多困难而又深刻的数学定理的来源。
